\scanpage{22}
\uline{Remarque 18} Il y a lieu d'appeler un foncteur \inserted{(additif ?)} satisfaisant les conditions du th. 12, sauf l'exactitude à gauche (ii), un \uline{foncteur constructible}, Il faudrait \deleted{(pour $X$ sur $Y$)} \unsure{définir|dégager|disposer} de \unsure{m\textsuperscript{e}|même} la notion de \unsure{foncteur|foncteurs} \unsure{constructible|constructibles} $\operatorname{Qcoh}(Y)\longrightarrow\operatorname{Qcoh}(X)$. Il y aurait un \unsure{sorite|suite} : composé de foncteurs constructibles est \unsure{itou|item}, foncteurs $\underline{\operatorname{Ext}}^i,\underline{\operatorname{Tor}}_j,R^{\unsure{n|q}} f_*$ sont constructibles \ldots{}...

\begin{center}\uline{Applications}\end{center}

\uline{Théorème} 19 Soit $f:X\longrightarrow S$ morphisme propre, de prés. finie et plat. Alors pour tt $Z$ \uuline{affine} sur $X$, $f_*(Z)=\displaystyle\prod_{X/S} Z/X$ est représentable par un schéma affine sur $S$. Si $Z/X$ est de prés. finie (resp de type fini) alors il en est de \unsure{m\textsuperscript{e}|même} de $\displaystyle\prod_{X/S}Z/X$ sur $S$.

\uline{Dém} On peut supposer $S$ affine, \deleted{soit} \inserted{donc $X$ qu. cpt \unsure{et|+} \inserted{(\unsure{qu|qu.})} séparé} \deleted{\unsure{$Z_0\quad Z$}} Alors $Z=\underset{\longleftarrow}{\lim}\,Z_i$ \inserted{($\underset{\longleftarrow}{\lim}$ filtrante)}, $Z_i$ affine de présentation finie sur $X$, d'où $f_*(Z)=\underset{\longleftarrow}{\lim}\,f_*(Z_i)$, et il suffit de prouver pour le premier point que les $f_*(Z_i)$ sont repr. par schémas affines de prés. finie sur $S$. Pour le deuxième point, il suffira alors de traiter le cas $Z$ de t.f. sur $X$.
% archive sheet number: 21
